% =========================================================
% Proof: Constant Multiple Rule (Global Form)
% Source: volume-iii/analysis/differentiation/notes/notes-algebra.tex
% Mode: standard
% =========================================================

\subsection*{Constant Multiple Rule — Global Form}
\label{prf:constant-multiple-rule-global}

\begin{remark*}[Return]
\hyperref[thm:constant-multiple-rule]{$\leftarrow$ Back to Theorem (Constant Multiple Rule) in Notes}
\end{remark*}

\begin{proof}
Let $I \subseteq \mathbb{R}$ be an open interval, let $f : I \to \mathbb{R}$
be differentiable on $I$, and let $\alpha \in \mathbb{R}$. Set $h := \alpha f$.
The claim is that $h$ is differentiable on $I$ and $h' = \alpha f'$.

Let $c \in I$ be arbitrary. Since $f$ is differentiable at $c$, by
\ByThm{Constant Multiple Rule}, $h$ is differentiable at $c$ and
$h'(c) = \alpha f'(c)$. Since $c \in I$ was arbitrary, $h$ is differentiable
at every point of $I$ and the derivative functions satisfy $h' = \alpha f'$
on $I$.
\end{proof}

\begin{remark*}[Proof shape]
Universal instantiation of the pointwise Constant Multiple Rule at an
arbitrary $c \in I$. No new technique — the global statement is the pointwise
statement read as a function identity on $I$.

\FlashProof{Fix arbitrary $c \in I$. Apply the pointwise Constant Multiple
  Rule to get $h'(c) = \alpha f'(c)$. Since $c$ was arbitrary, $h' = \alpha
  f'$ on $I$.}
\end{remark*}

\begin{remark*}[Dependencies]
\begin{itemize}
  \item \hyperref[thm:constant-multiple-rule]{Constant Multiple Rule}:
    The pointwise result instantiated at each $c \in I$.
\end{itemize}

\FlashUses{thm:constant-multiple-rule}
\end{remark*}

% =========================================================
% Proof: Sum Rule (Global Form)
% Source: volume-iii/analysis/differentiation/notes/notes-algebra.tex
% Mode: standard
% =========================================================

\subsection*{Sum Rule — Global Form}
\label{prf:sum-rule-global}

\begin{remark*}[Return]
\hyperref[thm:sum-rule]{$\leftarrow$ Back to Theorem (Sum Rule) in Notes}
\end{remark*}

\begin{proof}
Let $I \subseteq \mathbb{R}$ be an open interval and let $f, g : I \to \mathbb{R}$
be differentiable on $I$. Set $h := f + g$. The claim is that $h$ is
differentiable on $I$ and $h' = f' + g'$.

Let $c \in I$ be arbitrary. Since $f$ and $g$ are differentiable at $c$, by
\ByThm{Sum Rule}, $h$ is differentiable at $c$ and $h'(c) = f'(c) + g'(c)$.
Since $c \in I$ was arbitrary, $h$ is differentiable at every point of $I$
and $h' = f' + g'$ on $I$.
\end{proof}

\begin{remark*}[Proof shape]
Universal instantiation of the pointwise Sum Rule at an arbitrary $c \in I$.

\FlashProof{Fix arbitrary $c \in I$. Apply the pointwise Sum Rule to get
  $h'(c) = f'(c) + g'(c)$. Since $c$ was arbitrary, $h' = f' + g'$ on $I$.}
\end{remark*}

\begin{remark*}[Dependencies]
\begin{itemize}
  \item \hyperref[thm:sum-rule]{Sum Rule}: The pointwise result instantiated
    at each $c \in I$.
\end{itemize}

\FlashUses{thm:sum-rule}
\end{remark*}

% =========================================================
% Proof: Product Rule (Global Form)
% Source: volume-iii/analysis/differentiation/notes/notes-algebra.tex
% Mode: standard
% =========================================================

\subsection*{Product Rule — Global Form}
\label{prf:product-rule-global}

\begin{remark*}[Return]
\hyperref[thm:product-rule]{$\leftarrow$ Back to Theorem (Product Rule) in Notes}
\end{remark*}

\begin{proof}
Let $I \subseteq \mathbb{R}$ be an open interval and let $f, g : I \to \mathbb{R}$
be differentiable on $I$. Set $h := fg$. The claim is that $h$ is
differentiable on $I$ and $h' = f'g + fg'$.

Let $c \in I$ be arbitrary. Since $f$ and $g$ are differentiable at $c$, by
\ByThm{Product Rule}, $h$ is differentiable at $c$ and
$h'(c) = f'(c)g(c) + f(c)g'(c)$. Since $c \in I$ was arbitrary, $h$ is
differentiable at every point of $I$ and $h' = f'g + fg'$ on $I$.
\end{proof}

\begin{remark*}[Proof shape]
Universal instantiation of the pointwise Product Rule at an arbitrary $c \in I$.

\FlashProof{Fix arbitrary $c \in I$. Apply the pointwise Product Rule to get
  $h'(c) = f'(c)g(c) + f(c)g'(c)$. Since $c$ was arbitrary,
  $h' = f'g + fg'$ on $I$.}
\end{remark*}

\begin{remark*}[Dependencies]
\begin{itemize}
  \item \hyperref[thm:product-rule]{Product Rule}: The pointwise result
    instantiated at each $c \in I$.
\end{itemize}

\FlashUses{thm:product-rule}
\end{remark*}

% =========================================================
% Proof: Quotient Rule (Global Form)
% Source: volume-iii/analysis/differentiation/notes/notes-algebra.tex
% Mode: standard
% =========================================================

\subsection*{Quotient Rule — Global Form}
\label{prf:quotient-rule-global}

\begin{remark*}[Return]
\hyperref[thm:quotient-rule]{$\leftarrow$ Back to Theorem (Quotient Rule) in Notes}
\end{remark*}

\begin{proof}
Let $I \subseteq \mathbb{R}$ be an open interval, let $f, g : I \to \mathbb{R}$
be differentiable on $I$ with $g(x) \neq 0$ for all $x \in I$, and set
$h := f/g$. The claim is that $h$ is differentiable on $I$ and
$h' = (f'g - fg')/g^2$.

Let $c \in I$ be arbitrary. Since $g(c) \neq 0$, by \ByThm{Quotient Rule},
$h$ is differentiable at $c$ and
\[
  h'(c) = \frac{f'(c)g(c) - f(c)g'(c)}{(g(c))^2}.
\]
Since $c \in I$ was arbitrary and $g(x) \neq 0$ on all of $I$, $h$ is
differentiable at every point of $I$ and $h' = (f'g - fg')/g^2$ on $I$.
\end{proof>

\begin{remark*}[Proof shape]
Universal instantiation of the pointwise Quotient Rule at an arbitrary
$c \in I$. The global hypothesis $g(x) \neq 0$ on all of $I$ ensures the
pointwise condition $g(c) \neq 0$ is satisfied at every application.

\FlashProof{Fix arbitrary $c \in I$. Since $g(c) \neq 0$, apply the pointwise
  Quotient Rule to get $h'(c) = (f'(c)g(c) - f(c)g'(c))/(g(c))^2$. Since
  $c$ was arbitrary, $h' = (f'g - fg')/g^2$ on $I$.}
\end{remark*}

\begin{remark*}[Dependencies]
\begin{itemize}
  \item \hyperref[thm:quotient-rule]{Quotient Rule}: The pointwise result
    instantiated at each $c \in I$, valid since $g(c) \neq 0$ for all
    $c \in I$.
\end{itemize}

\FlashUses{thm:quotient-rule}
\end{remark*}